\section{研究目标}

结合研究点一与研究点二的研究内容，本研究需要证明或达成以下效果：

\begin{enumerate}
  \item 在与AgentDyn运行时结构对应的调用前中介点上，构建版本化任务契约与确定性动作门控机制，能够将逐次工具调用的授权闭包收敛为\texttt{AUTHORIZED}/\texttt{UNRESOLVED}/\texttt{UNAUTHORIZED}三态判定，并把判定结果映射为可执行的五级防御动作；
  \item 在共享拒绝机会、相同调用前状态与冻结预算的受控条件下，证明面向规划器可见的拒绝反馈披露粒度对同授权安全完成率、未授权提交发生率与预算内攻击自适应成功率存在可观测的条件效应；
  \item 在与AgentDyn运行时事件对应的证据采集点上，构建区分硬边与软边的增量执行证据图，并证明基于两条在线强制路径规则的检测机制能够识别跨步骤传播且已闭合的高危路径；
  \item 证明面向可执行阻断原语的有界字典序最小覆盖方法，能够在覆盖已识别路径证据包的前提下，较逐路径确定性阻断策略产生更小的干预范围与更高的任务效用保持；
  \item 在独立评测视角下，证明运行轨迹图防御方法输出的风险子图与根因节点定位相对无图结构基线具有更高的注入点包含召回率与根因定位准确率；
  \item 完成两个研究点在统一原型系统上的集成，并在AgentDyn基准上产出可复现的主实验、消融实验与统计检验结果，支撑上述各项效果的实证结论。
\end{enumerate}
